%%%%%%%%%%%%%%%%%%%%%%%%%%%%%%%%%
%
%       EXERCÍCIO
%
%%%%%%%%%%%%%%%%%%%%%%%%%%%%%%%%%

\ifdefstring{\atividade}{prova}{%
    \ifdefstring{\modo}{objetivo}{%
        \renewcommand{\valorquestao}{\ValorQObj\ ponto}
    }{%
        \renewcommand{\valorquestao}{\ValorQDisc\ pontos}
    }%
}%

\begin{exercicioBanco}[\valorquestao]
Em uma turma, \(50\) estudantes responderam a uma pesquisa sobre duas redes sociais, \(A\) e \(B\).
Sabe-se que:

\begin{enumerate}[\(\bullet\)]
    \begin{multicols}{2}
        \item \(24\) estudantes utilizam a rede social \(A\);
        \item \(10\) estudantes utilizam ambas as redes sociais;
        \item \(16\) estudantes não utilizam nenhuma das duas redes.
    \end{multicols}
\end{enumerate}

Sabendo que todos os \(50\) estudantes responderam à pesquisa, o número de estudantes que utilizam a rede social \(B\) é:

% Define as alternativas
\newcommand{\alternativas}{%
    \begin{center}
        \begin{tabularx}{\textwidth}{XXXXX}
            (a) \(18\). &
            (b) \(20\). &
            (c) \(22\). &
            (d) \(24\). &
            (e) \(26\).
        \end{tabularx}
    \end{center}
}

% Define a resposta correta
\newcommand{\resposta}{B}

% Lógica condicional para exibição
\ifdefstring{\atividade}{lista}{%
    \alternativas
    \vspace{0.5em}
    
    \noindent\textbf{Resposta:} letra \textbf{\resposta}.
}{%
    \ifdefstring{\modo}{objetivo}{%
        \alternativas
    }{%
        % modo = discursiva → não mostra alternativas
    }
}
\end{exercicioBanco}

